% 纤维丛结构图 (Fiber Bundle Structure)
% 描述: 展示纤维丛 E → B 的投影映射，局部平凡化
% 数学概念: 纤维丛、底空间、全空间、纤维、局部平凡化
\documentclass[border=10pt]{standalone}
\usepackage{tikz}
\usepackage{amsmath,amssymb}
\usetikzlibrary{arrows.meta,positioning,calc,patterns,3d}

\begin{document}
\begin{tikzpicture}[
    >=Stealth,
    scale=1.0,
    base/.style={
        draw=blue!60!black,
        fill=blue!15,
        line width=1pt,
        rounded corners=3pt
    },
    fiber/.style={
        draw=red!70!black,
        fill=red!20,
        line width=0.8pt
    },
    total/.style={
        draw=purple!70!black,
        fill=purple!10,
        line width=1pt,
        opacity=0.8
    },
    projection/.style={
        ->,
        line width=1.5pt,
        color=orange!80!black,
        dashed
    },
    section/.style={
        line width=1.2pt,
        color=green!60!black,
        smooth
    }
]

% ============= 底空间 B =============
% 绘制底空间 (曲面表示)
\draw[base] 
    (0,0) .. controls (2,1) and (4,0.5) .. (6,0)
          .. controls (6.5,-0.5) and (6,-1) .. (5,-1.2)
          .. controls (3,-1.5) and (1,-1) .. (0,0);
\node[font=\bfseries, color=blue!60!black] at (3, -1.7) {底空间 $B$};

% 底空间上的点
\coordinate (b) at (3, -0.3);
\fill[black] (b) circle (2pt);
\node[below] at (b) {$b$};

% ============= 全空间 E =============
% 在全空间区域绘制纤维丛
\begin{scope}[yshift=2.5cm]
    % 全空间外轮廓
    \draw[total] 
        (-0.5,0) .. controls (1.5,2.5) and (4.5,2) .. (6.5,0)
              .. controls (7,-1.5) and (5.5,-2.5) .. (4,-2)
              .. controls (2,-1.5) and (0.5,-1.5) .. (-0.5,0);
    
    % 多条纤维 (垂直线段)
    \foreach \x in {0.5, 1.5, 2.5, 3.5, 4.5, 5.5} {
        \draw[fiber] (\x, -1.2) -- (\x, 1.5);
    }
    
    % 特殊纤维 F_b (在点b上方)
    \coordinate (fb_top) at (3, 1.5);
    \coordinate (fb_bottom) at (3, -1.2);
    \draw[fiber, line width=2pt, color=red!80!black] (fb_bottom) -- (fb_top);
    \node[color=red!80!black, right] at (3.1, 0.2) {$F_b = \pi^{-1}(b)$};
    
    % 全空间中的点
    \coordinate (e) at (3, 0.5);
    \fill[black] (e) circle (2pt);
    \node[above right] at (e) {$e$};
    
    % 全空间标签
    \node[font=\bfseries, color=purple!70!black] at (5.5, 1.8) {全空间 $E$};
    
    % 截面
    \draw[section] 
        (0.8, -0.3) .. controls (2, 0.5) and (3, 0.2) .. (4, 0.4) 
                    .. controls (5, 0.6) .. (5.8, 0.1);
    \node[color=green!60!black, above] at (5.5, 0.3) {$\sigma$};
\end{scope}

% ============= 投影映射 =============
\draw[projection] (e) -- (b);
\node[color=orange!80!black, right] at (3.3, 1.2) {$\pi$};

% ============= 局部平凡化图示 =============
\begin{scope}[xshift=8cm, yshift=0cm]
    % 标题
    \node[font=\bfseries] at (2, 4) {局部平凡化};
    
    % 局部邻域 U
    \draw[base, fill=blue!10] (0, 0) ellipse (1.5 and 1);
    \fill[black] (0, 0) circle (1.5pt);
    \node[below] at (0, -0.1) {$b$};
    \node[color=blue!60!black] at (0, -1.4) {$U \subset B$};
    
    % 映射箭头
    \draw[->, line width=1.2pt, color=orange!80!black] (1.8, 2.5) -- (2.8, 2.5);
    \node[color=orange!80!black, above] at (2.3, 2.5) {$\varphi$};
    
    % 乘积空间 U × F
    \begin{scope}[xshift=4cm, yshift=0.5cm]
        % 矩形背景
        \fill[gray!15, rounded corners=3pt] (-1.5, -1.2) rectangle (1.5, 2.2);
        \draw[gray!50, step=0.3] (-1.5, -1.2) grid (1.5, 2.2);
        
        % 底边 U
        \draw[base, line width=2pt] (-1.2, -0.8) -- (1.2, -0.8);
        \node[color=blue!60!black, below] at (0, -0.9) {$U$};
        
        % 垂直纤维
        \foreach \x in {-0.8, -0.4, 0, 0.4, 0.8} {
            \draw[fiber, line width=0.6pt] (\x, -0.8) -- (\x, 1.5);
        }
        
        % 标签
        \node[color=red!70!black, right] at (1, 0.5) {$F$};
        \node at (0, 2.5) {$U \times F$};
        
        % 对应点
        \fill[black] (0, 0.3) circle (1.5pt);
        \node[right, font=\small] at (0.1, 0.3) {$(b, f)$};
    \end{scope}
    
    % 说明
    \node[anchor=north, font=\scriptsize, text width=6cm, align=center] 
        at (2, -2) {
        局部平凡化: $\varphi: \pi^{-1}(U) \xrightarrow{\cong} U \times F$\\
        使得纤维丛局部看起来像乘积空间
    };
\end{scope}

% ============= 整体结构公式 =============
\node[anchor=south, font=\small, draw=gray!50, fill=white, rounded corners, inner sep=8pt] 
    at (4, -3.5) {
    \begin{tabular}{c}
        \textbf{纤维丛结构:} $(E, B, \pi, F)$ \\[3pt]
        $\pi: E \to B$ 是连续满射，每点 $b \in B$ 有邻域 $U$ 使得 
        $\pi^{-1}(U) \cong U \times F$ \\[3pt]
        \textbf{截面:} $\sigma: B \to E$ 满足 $\pi \circ \sigma = \text{id}_B$
    \end{tabular}
};

\end{tikzpicture}
\end{document}
